\zhabstract

本文档使用UESTC Thesis模板编译生成，用于展示电子科技大学学位论文的基本结构与排版效果。

本文档涵盖中英文摘要、章节层次、数学公式、图表、文献引用、主要符号表、缩略词表、附录及
攻读学位期间取得的成果等常用内容。文中的文字、数据和图表仅供排版演示，使用者可根据实际
论文内容进行修改或替换。

不同学科的学位论文在内容组织和写作方式上有所不同。使用者可通过中国知网、万方数据等学术
数据库查阅本学科的优秀学位论文，了解常见的内容组织与表达方式。论文格式应以电子科技大学
研究生院发布的学位论文撰写规范及相关模板为准，相关文件可从以下网址获取：
\url{https://gr.uestc.edu.cn/xiazai/114/3917}。本模板依据学校最新规范及官方Word模板设计，
对其主要结构与版式进行了高精度复刻，使用者可放心用于论文写作与排版；正式提交前，仍应
结合所在学院的具体要求进行最终核对。

\keywords{电子科技大学；学位论文；LaTeX模板；论文排版}

\enabstract

This document is compiled using the UESTC Thesis template to demonstrate the
basic structure and typesetting of dissertations and theses at the University
of Electronic Science and Technology of China.

It covers Chinese and English abstracts, heading levels, mathematical equations,
figures, tables, citations, lists of symbols and abbreviations, appendices, and
academic achievements during the degree program. The text, data, and figures are
provided solely for typesetting demonstration and may be modified or replaced
according to the user's thesis.

The organization and writing conventions of dissertations and theses vary across
disciplines. Users may consult outstanding dissertations and theses in their
fields through academic databases such as CNKI and Wanfang Data to learn common
approaches to organization and scholarly expression. Thesis formatting should
follow the latest writing standards and related templates published by the UESTC
Graduate School, available at \url{https://gr.uestc.edu.cn/xiazai/114/3917}.
Based on these official documents and Word templates, the UESTC Thesis template
reproduces their principal structure and layout with high fidelity and can be
used confidently for thesis writing and typesetting. Before final submission,
users should nevertheless confirm any additional requirements issued by their
school or department.

\enkeywords{UESTC; Dissertations and Theses; LaTeX Template; Thesis Typesetting}
